\documentclass[12pt]{article}
\usepackage[utf8]{inputenc}
\usepackage[T1]{fontenc}
\usepackage{hyperref}
\title{Sidetic Language Resources for K-12 Curriculum Development}
\author{Compiled by Claude Sonnet 5.5 (claude-sonnet-5-5) for Andr\'as Kornai}
\date{}
\begin{document}
\maketitle

\section{Introduction}
Sidetic is an extinct Anatolian (Indo-European) language of Side on the Pamphylian coast. It is known only from coin legends, roughly 5th to 3rd century BCE, and a small number of short inscriptions. It has no living speakers and no native-speaker community. A K-12 curriculum taught in Sidetic is therefore not possible. At most, Sidetic could be a topic within history, linguistics or classics units. The search found no school materials.

\section{Primary Educational Resources}
None were found. The search turned up no primers, textbooks or school programmes for Sidetic.

\section{Linguistic and Reference Materials}
\begin{itemize}
\item Wikipedia, ``Sidetic language'': overview of the corpus, the script and its decipherment. \url{https://en.wikipedia.org/wiki/Sidetic_language}
\item Britannica, ``Sidetic language'': describes it as very sparsely documented. \url{https://www.britannica.com/topic/Sidetic-language}
\item F. Woudhuizen, ``On the reading and interpretation of the two longer Sidetic inscriptions S I.2.1 and S I.2.5'' (a PDF on antiquitasviva.com). Woudhuizen's interpretation is contested. \url{https://antiquitasviva.com/wp-content/uploads/2021/07/70.1-2.04.-woudhuizen-f.-on-the-reading-and-interpretation-of-the-two-longer-sidetic-inscriptions-s-i.2.1-and-s-i.2.5.pdf}
\item University of Verona repository record on the new Sidetic inscriptions found at Side since 2020. \url{https://iris.univr.it/handle/11562/1164487}
\item Hurriyet Daily News, news report on research that raised the known alphabet to 31 letters. \url{https://www.hurriyetdailynews.com/amp/researchers-identify-31-letters-in-ancient-sidetic-language-222957}
\end{itemize}

\section{Religious and Community Resources}
None were found. There is no living community. The inscriptions are mostly votive texts.

\section{Digital and Audio Resources}
No audio, software or corpus site was found. The only related item seen is a numismatic auction page showing a stater from Side with its legend: \url{https://leunumismatik.com/de/lot/18/238}.

\section{Curriculum Development Context}
The language is extinct, the corpus is tiny and it is not fully deciphered. Any curriculum use would be enrichment in history, archaeology or writing-systems units, using the sources above. The 31-letter alphabet could serve as an example of an ancient Anatolian script. No Sidetic-medium instruction is feasible.
\end{document}